\section{案例 1：学习启动失败}

\subsection*{情境}
一名学生午饭后本打算阅读一个较难的章节，却转而打开短视频信息流，连续刷了四十分钟。

\subsection*{框架映射}
\begin{itemize}[nosep]
  \item \textbf{当前状态：} 午饭后的被动新奇消费。
  \item \textbf{目标状态：} 坐下后聚焦证明的阅读，并产生活跃笔记输出。
  \item \textbf{主导变量：} 高 $\Reward$、高 $\Switch$，以及高 $\Env$。
  \item \textbf{错过的决策窗口：} 饭后立刻重新进入的那段时间，在手机把旧机制重新搭起来之前。
\end{itemize}

\subsection*{机制解释}
这个案例最适合被理解为一次局部状态转变失败。当前状态带有一种很容易再次进入的类吸引子特征：低努力、快速新奇，以及熟悉的姿势与工具配置。目标状态之所以没有出现，并不只是因为学生缺少价值认同；更关键的是，它本身稳定性弱，而且事前准备不足。涌现章节又提供了第二层解释：如果午饭后反复拿起手机，这条时间边界就会形成路径依赖，于是同一个时间点不再是中性的，而会变成一个背负历史的提示线索。

\subsection*{证据车道纪律}
这里，强经验意识研究来源只支持一种有限的迁移：有组织的状态及其转变，确实可能表现出结构化的稳定性与动态变化。因此，这个案例使用的是\textbf{从强到中等}的翻译，而不是直接证明。关于睡眠转变分岔的论文，可以作为“为什么转变时机重要”的模型，但它并不能证明每一次午休结束都在严格神经科学意义上构成一个分岔点。亚稳态语言在\textbf{中等}意义上是有用的：工作状态必须足够稳定，才能被保持住；同时，中断后还得足够容易恢复。

\subsection*{操作性修复}
\begin{enumerate}[nosep]
  \item 在吃饭前写下下一段证明的精确目标。
  \item 让笔记本停留在尚未完成的那一行。
  \item 把手机移出眼前的工作平面。
  \item 把休息后的第一个动作定义为“坐下，读最后两行，再补写一句延续句子”。
\end{enumerate}

\subsection*{迁移教训}
学习启动失败通常不是靠提高道德压力来解决的。它更常见的修复方式，是让预期的重新进入轨迹更短、更可见，并且更少依赖回到桌前那一刻的临时设计。